% Bài làm buổi 20/10 — Lúa
\begin{problem}[9DS-BPT-004]
$\dfrac{x - 2}{4} - \dfrac{2x + 1}{3} \geq \dfrac{1}{2}$
Nhân hai vế với $12$: $3(x - 2) - 4(2x + 1) \geq 6$
$3x - 6 - 8x - 4 \geq 6$
$-5x - 10 \geq 6$
$-5x \geq 16$
$x \leq -\dfrac{16}{5}$
Vậy $x \leq -\dfrac{16}{5}$.
\end{problem}

\begin{problem}[9DS-CAN-004]
a) $A = 2\sqrt{4 \cdot 3} - 3\sqrt{9 \cdot 3} + \sqrt{25 \cdot 3} + \dfrac{1}{2}\sqrt{16 \cdot 3}$
$A = 4\sqrt{3} - 9\sqrt{3} + 5\sqrt{3} + 2\sqrt{3} = 2\sqrt{3}$
b) $\sqrt{18a^2} = \sqrt{9 \cdot 2 \cdot a^2} = 3a\sqrt{2}$
c) $5\sqrt{2} = \sqrt{50}$, $4\sqrt{3} = \sqrt{48}$, $3\sqrt{6} = \sqrt{54}$, $2\sqrt{13} = \sqrt{52}$
Vì $48 < 50 < 52 < 54$ nên $4\sqrt{3} < 5\sqrt{2} < 2\sqrt{13} < 3\sqrt{6}$
\end{problem}

\begin{problem}[9DS-PT2-005]
a) $\Delta = (-5)^2 - 4 \cdot 1 \cdot 3 = 13 > 0$ nên phương trình có hai nghiệm phân biệt.
b) Theo định lý Viète: $x_1 + x_2 = 5$, $x_1 x_2 = 3$
$A = x_1^2 + x_2^2 = (x_1 + x_2)^2 - 2x_1 x_2 = 25 - 6 = 19$
$B = \dfrac{1}{x_1} + \dfrac{1}{x_2} = \dfrac{x_1 + x_2}{x_1 x_2} = \dfrac{5}{3}$
$C = x_1^3 + x_2^3 = (x_1 + x_2)^3 - 3x_1 x_2(x_1 + x_2) = 125 - 45 = 80$
\end{problem}

\begin{problem}[9DS-HSB2-006]
a) $(P)$ đi qua $M(2; -2)$ nên $-2 = a \cdot 2^2$, suy ra $a = -\dfrac{1}{2}$
Vì $a < 0$ nên $(P)$ nằm phía dưới trục hoành.
b) Phương trình hoành độ giao điểm: $-\dfrac{1}{2}x^2 = x - 4$
$x^2 + 2x - 8 = 0$
$(x - 2)(x + 4) = 0$
$x = 2$ hoặc $x = -4$
Với $x = 2$ thì $y = -2$; với $x = -4$ thì $y = -8$.
Vậy các giao điểm là $(2, -2)$ và $(-4, -8)$.
\end{problem}
